\documentclass[../proyecto.tex]{subfiles}
\graphicspath{{\subfix{../images/}}}
\begin{document}

\chapter{Investigación propuesta}

% máximo diez páginas

\section{Presentación del proyecto}

\section{Lugar de enunciación}

Mi práctica artística y académica es una mezcla de confluencias y contradicciones.

Soy una persona chilena, nacida a fines del siglo veinte, que le encanta la computación pero no los videojuegos.

Partí 

\section{Motivaciones}

Creo en la importancia de que existan nuevos cánones sobre luthería electrónica, que existan más miradas sobre lo que implica sintetizar sintetizadores.

\section{Marco teórico / definiciones conceptuales}

En industria y en español existe muchas veces una concepción de sintetizadores y teclados, pero en esta tesis nos concentraremos en otros tipos de interfaces \footcite{bjorn_push_2021}.

\section{Problema / objeto de estudio / corpus}


\section{Criterios de selección del corpus}

En este corpus se incluyen tanto libros de la última década que sistematizan prácticas artísticas y pedagógicas \footcite{pearson_electronic_2024}.

\section{Referencias críticas sobre objeto de estudio / estado del arte}

\section{Referentes artísticos y contexto en que se sitúa la práctica artística}

Peter Blasser y sus líneas de sintetizadores Ciat-Lonbarde, y Shbobo.

La empresa japonesa KORG y su línea de sintetizadores programables.

\section{Novedad de la propuesta}



\end{document}